\providecommand{\PushA}{\frac{5422587}{6103515625}}
\providecommand{\PushK}{890666}
\section{Large primes \texorpdfstring{$p\ge2^{512}$}{p >= 2\textasciicircum 512} and the external index}\label{sec:tail-index}

\paragraph{What this section proves and how to read it.}
Large primes (called the \emph{tail} below) are one atom per prime,
of both residue classes; this is the only family whose size grows
with $N$. The section proves (CP1) and summability with prices
$O(p^{-3/2+2a})$, and evaluates the conductor factor $R_I$ and the
terms $J,W$ at the baseline choice $\Kint=\PushK$. ``Tail'' here always means large
primes; depth tails $\delta_{p,j}$ are a different notion
(\cref{sec:comparison-theorem}).
Two parts of the construction are not blocks of field primes of bounded
size. The \emph{tail} consists of all primes $p\ge\Pfour=2^{512}$, of both
residue classes modulo $4$, one atom $\{p\}$ per prime with $e_p=1$; it is
the only family of atoms whose number grows with $N$. The tail law is the
reference of the tail datum (\cref{def:record:upper-bands}(e)) at a single
prime; it is not the law of the last band, whose data are class-three data.
The \emph{small modulus} $\Mzero$ (formula in
\cref{subsec:index-law}) is not a coordinate of $G_N$; it enters
only through the external index of the six queries. We show that the tail law satisfies
(CP1) at every tail prime (\cref{subsec:tail-law}) and \ref{it:record:R4} at
the primes $p\ge\Pfive=2^{65536}$, as \cref{lem:record:prefix-conductor}
requires (\cref{subsec:tail-mixing}), and that the tail comparison prices are
$O(p^{-3/2+2a})$ per prime, so that they enter $S$ and $A$ as a convergent
series, satisfy the summability (CP8)(c), and make the product $\Pi_{\rm H}$
of \cref{lem:record:density-lp} finite (\cref{subsec:tail-prices}). For the
index we evaluate the conductor moment $R_I$ and the terms $J$ and $W$ of
\cref{thm:comparison} at the baseline choice $\Kint=\PushK$, with the finite prime-moment bound
of \cref{lem:index-moments} (\cref{subsec:index-law}). These enter the closing
(\cref{sec:closing}).

\subsection{The reference law at one large prime}\label{subsec:tail-law}

The upper reference is defined band by band with class-three data
(\cref{def:record:upper-bands}). Beyond $\Pfour$ every prime is an atom of
its own and carries, at the single prime, the reference of the \emph{tail
datum} (\cref{def:record:upper-bands}(e)): the inputs and weights recorded for
the range $[2^{64},2^{512})$ in
\path{research/roadmap_2026_10_18/item12/upper_v2_local.json}, whose inputs take
the larger value over the two classes. The square-clique prefix ends at $\Pfour$
(\cref{lem:frontier:iidlocal}),
so the tail also contains the primes $p\equiv1\pmod4$. This is the only place
where a two-layer law is used at class-one primes, and the only place where
the class-one cases of the spectral criterion are needed
(\cref{subsec:advance:parity}).

\begin{definition}[The tail law]\label{def:tail-law}
Let $p\ge\Pfour$ be prime and $r=p^{-1/2}$. The \emph{tail law} at $p$ is the
reference $\mu_p=\widetilde\mu_p/P_\mu(r)$ of \cref{def:upper-reference} for
the single prime $\Blk=\{p\}$, with the weights $w_{\mathfrak t},v_{\mathfrak t}$
of the tail datum (\cref{def:record:upper-bands}(e)):
\[
 \widetilde\mu_p=\sum_{\mathfrak t}w_{\mathfrak t}r^{\deg\mathfrak t}\pi_{\mathfrak t,2}
 +r\sum_{\mathfrak t}v_{\mathfrak t}r^{\deg\mathfrak t}
 \operatorname{Diag}_\iota(\pi_{\mathfrak t,1}),\qquad
 P_\mu(r)=\sum_{d=0}^7c_dr^d,
\]
with the components $\pi_{\mathfrak t,j}$ of \cref{subsec:layer72:upper} at
$p$. Here $c_0=w_{\rm LL}=1$ is the weight of the \emph{root} $\pi_{\rm LL,2}$,
the uniform law on the injective strict-square labellings of the two-layer root
graph, and $c_d$ is the total weight of the records of overall degree $d$
(\eqref{eq:independent-fold-polynomial}):
\begin{gather*}
 c_0=1,\quad c_1=\tfrac{5988}{5},\quad c_2=\tfrac{15232626}{25},\quad
 c_3=\tfrac{24807984744}{175},\\
 c_4=\tfrac{50638290432}{5},\quad c_5=\tfrac{621170191872}{5},\quad
 c_6=\tfrac{1877429977088}{7},\quad c_7=\tfrac{970455252992}{7}.
\end{gather*}
The \emph{tail coefficients} are
\begin{equation}\label{eq:tail-coefficients}
 C_j=\sum_{d\ge j}c_d\Pfour^{-(d-j)/2}\quad(j=1,2,3),\qquad
 S_{72}(p)=\frac{2^{92}}p+\frac{C_1}{p^{1/2}}+\frac{C_2}p
 +\frac{C_3}{p^{3/2}};
\end{equation}
$C_1<1197.601$, $C_2<609305.05$ and $C_3<1.4176\cdot10^8$ (item (S7), the $C_j$ bounds, of
\cref{cert:record:analytic-stage}, whose checker recomputes them exactly from
the $c_d$).
\end{definition}

\begin{lemma}[The tail reference]\label{lem:tail-law}
Let $p\ge\Pfour$ be a prime of either class modulo $4$, and $r=p^{-1/2}$.
\begin{enumerate}[label=(\roman*)]
\item The inputs of the tail datum (\cref{def:record:upper-bands}(e)) are
 valid on $\{p\}$ in the sense of \cref{subsec:upper-reference-rp}.
\item $\mu_p$ is a probability measure on $\F_p^V$, reflection positive at
 every cut of the family (the chamber cuts of both coordinates and
 $\iota$), invariant under the square-affine maps $z_v\mapsto sz_v+u$
 ($s\in\Qsq$, $u\in\F_p$), with uniform one-slot marginals, and admissible
 (\cref{def:admissible}); so the tail atoms satisfy
 \ref{it:record:R1}--\ref{it:record:R3}.
\item $\widetilde\mu_p=\pi_{\rm LL,2}+\widetilde\rho_p$ with $\widetilde\rho_p\ge0$
 of mass $h_p=P_\mu(r)-1\le C_1p^{-1/2}$. The root accepts the twelve
 queries \eqref{eq:upper-twelve}, hence $\mathcal E_6$, and each query pair
 is a cross-layer edge of the root graph.
\end{enumerate}
\end{lemma}

\begin{proof}
(i) The inputs of the tail datum are integer ceilings of constants evaluated
at $p_0=2^{64}$ with the block length $L_{\rm u}=50000/27$, that is, for blocks
of at most $L_{\rm u}/\log2^{64}<42$ primes, and they cover both prime
classes. The nonroot caps are the largest of four exact tests, two classes
and two Fourier signs (\cref{subsec:upper-nonroot},
\cref{cert:upper-band-inputs}(d)). The root caps $113$ (LL20) and $45$ (LL22)
come from the tests of both classes, the class-one ones with the class-one
Paley terms and the parity transfer (\cref{lem:upper-completed-root} with
$\mathfrak c\in\{1,3\}$, \cref{cert:upper-completed-root},
\cref{lem:spec2:paley-one,lem:spec2:parity-transfer}); at $2^{64}$ the
class-one tests give the recorded value in every cut except FE22 and EE22,
where the class-one and class-three test values are equal. The row inputs are the larger of the values of the
two classes (\cref{subsec:advance:product-rows}). By
\cref{lem:record:domain}(iii) for \ref{it:record:upper-F1}--\ref{it:record:upper-F6},
with \ref{it:record:upper-F5} in its Schur form, and by
\cref{lem:upper-completed-root} for the two root cuts (its threshold
hypotheses at $p_0$ are asserted by \cref{cert:upper-completed-root}), every
per-prime constant at a prime $p\ge2^{64}$ of either class is at most its
value at $2^{64}$; the rows are those of \ref{it:record:upper-F5}(b), with
$t_Y=1/50$ and $M=21/20$. The set $\{p\}$ is a block of one prime, so no tensor-product factor occurs and its block bounds are the per-prime bounds at $p$:
the normalized cut operators at $p$ are at least $-g\,r$ with the recorded
per-prime cap $g$, which is at most the inflated cap
$g\,e^{B_{\rm top}L_{\rm u}/(X\log X)}$ at $X=2^{64}$ (notation of \ref{it:record:upper-F6}) recorded as the input of the tail
datum, and the row data at $p$ satisfy \ref{it:upper-valid-rows} with
coefficients at most the row inputs. These are the conditions
\ref{it:upper-valid-root}--\ref{it:upper-valid-layer}.

(ii) By (i), \cref{lem:upper-reference-rp}(i), which uses nothing about the
set of primes beyond the validity of the inputs, applies to $\Blk=\{p\}$ with
the weights of the tail datum. These satisfy its $105$ inequalities exactly
(\cref{cert:upper-completion}, last paragraph), which do not involve $p$.
This gives (ii); the law of the atom $\{p\}$ does not depend on $N$.

(iii) $\pi_{\mathfrak t,j}$ has mass $m_{\mathfrak t}$ and $m_{\rm LL}=1$, so
$\widetilde\mu_p$ has mass $P_\mu(r)$; and $p^{-d/2}\le\Pfour^{-(d-1)/2}p^{-1/2}$
for $d\ge1$, which gives $h_p=\sum_{d\ge1}c_dp^{-d/2}\le C_1p^{-1/2}$. By
\cref{lem:upper-geometry}(a) every query
$(u_0,u_1,0)\to(\rot u_0,\rot u_1,1)$ of \eqref{eq:upper-twelve} is an edge of
the root graph joining its two layers, so the root, the uniform law on the
injective strict-square labellings of that graph, accepts the twelve queries;
the six queries of $\mathcal E_6$ are among them.
\end{proof}

\subsection{Mixing at the tail primes}\label{subsec:tail-mixing}

Hypothesis \ref{it:record:R4} is needed only at $p\ge\Pfive$: the atoms
below $\Pfive$, among them the tail primes in $[\Pfour,\Pfive)$, are finitely
many, their full orders divide the integer $Q_{\rm pre}$ of
\cref{def:record:functional}, and in \cref{lem:record:prefix-conductor} they
form the spectator prefix, which only has to be a contraction. The decay
comes from the root, whose conditional operators have norm $O(p^{-1/2})$;
the other records are an admixture of mass $O(p^{-1/2})$.

\begin{lemma}[Tail mixing]\label{lem:record:tail-mixing}
Let $p\ge\Pfour$ be prime, let $\mu_p$ be the tail law, and let
$\bar\nu_p=\nu_p/\nu_p(1)$ be its normalized accepted part, where $\nu_p$ is
the factor at $p$ of the comparison law (\cref{lem:tail-prices},
\cref{def:record:comparison-law}).
\begin{enumerate}[label=(\alph*)]
\item For every slot $v$, the conditional operators of $\mu_p$ and of
 $\bar\nu_p$ from slot $v$ to the other $71$ slots (\cref{def:mixing}) have
 squared norm at most $S_{72}(p)$ on mean-zero functions.
\item For each query $(u,v)\in\mathcal E_6$, the $(u,v)$-marginal of
 $\bar\nu_p$ is $\Unif(\Om)$, and its conditional operator from the pair
 $(u,v)$ to the other $70$ slots fixes constants, maps mean-zero functions
 to mean-zero functions and has squared norm at most $S_{72}(p)$ on them.
\item With $\gamma=255/1024$, $p^{2\gamma}S_{72}(p)$ is decreasing in $p$.
 Hence $S_{72}(p)<p^{-2\gamma}$ for every $p\ge\Pfive$, and
 \ref{it:record:R4} holds.
\end{enumerate}
\end{lemma}

\begin{proof}
\emph{The tail laws as mixtures.} By \cref{lem:tail-law}(iii),
$\mu_p=\varpi\pi_{\rm LL,2}+(1-\varpi)\widetilde\rho_p/h_p$ with
$\varpi=1/(1+h_p)$, and $0<h_p\le C_1p^{-1/2}\le S_{72}(p)-2^{92}/p$. The
comparison law keeps the root,
which accepts $\mathcal E_6$, keeps the other accepting records, replaces
each repaired record by a law of the same mass and deletes the others
(\cref{lem:tail-prices}); so $P_\mu(r)\nu_p=\pi_{\rm LL,2}+\rho_p$ with
$\rho_p\ge0$ of mass $0<h'_p\le h_p$ (the records of degree one are kept), and
$\bar\nu_p=\varpi'\pi_{\rm LL,2}+(1-\varpi')\rho_p/h'_p$ with
$\varpi'=1/(1+h'_p)$.

(a), (b). Apply \cref{lem:mixture} to these two-component mixtures, with $Z$
the labels of the slot $v$ or of the pair $(u,v)$. The components have the
same marginal on $Z$: one-slot marginals are uniform
(\cref{lem:tail-law}(ii), \cref{lem:tail-prices}(i)); the queried pair of
$\bar\nu_p$ is uniform on $\Om$ (\cref{lem:tail-prices}(i)), and so is that
of the root, which is square-affine invariant and accepts the query, the
square-affine maps acting transitively on $\Om$; hence so is that of
$\rho_p/h'_p$. The root graph has $72$ vertices and maximum degree $20$
(ten out- and ten in-neighbours at every vertex, \cref{lem:upper-geometry}(a)),
and the query pair is one of its edges, oriented from $u$ to $v$
(\cref{lem:tail-law}(iii)); by \cref{lem:trace-mixing}\ref{it:trace:72}
(\cref{cert:trace-mixing-72}), for $p\ge2^{512}$ the root has one-slot
conditional norm squared at most $2^{22}/p$ and pair conditional norm squared
at most $2^{92}/p$ on the mean-zero functions. The other component is a
contraction there. By \cref{lem:mixture} the conditional operators fix the
constants, map mean-zero functions to mean-zero functions, and have squared
norm at most $\varpi2^{92}/p+(1-\varpi)\le2^{92}/p+h_p\le S_{72}(p)$, and the
same with $\varpi'$ and $h'_p\le h_p$ for $\bar\nu_p$.

(c) $p^{2\gamma}S_{72}(p)=2^{92}p^{2\gamma-1}+C_1p^{2\gamma-1/2}
+C_2p^{2\gamma-1}+C_3p^{2\gamma-3/2}$ has nonnegative coefficients and
negative exponents, since $2\gamma=255/512<1/2$. At $p=\Pfive$,
$p^{2\gamma}=2^{32640}$, so (S7) gives $p^{2\gamma}S_{72}(p)<1$ for every
$p\ge\Pfive$. With (a) and (b), $\mu_p$ and $\bar\nu_p$ are
$(\gamma,v)$-mixing for every slot $v$, and $\bar\nu_p$ has the pair property
of \ref{it:record:R4}.
\end{proof}

The other inequality of (S7), $2^{16}S_{72}(2^{512})<6.7782\cdot10^{-70}$, would give
by the same monotonicity, since $\Pfour^{1/32}=2^{16}$, mixing with exponent
$1/64$ at every tail prime; it is not used.

\subsection{Comparison prices \texorpdfstring{$O(p^{-3/2+2a})$}{O(p\textasciicircum (-3/2+2a))} at large primes}\label{subsec:tail-prices}

At a tail prime the comparison law is built as at an upper block
(\cref{def:upper-comparison}), but for the twelve queries: a record is kept if
it accepts \eqref{eq:upper-twelve}. Its repair caps are the \emph{all-prime}
caps, since $p$ may be of class one. For each of the $94$ two-layer records
of degree two and three that do not accept the twelve queries
(\cref{sec:frontier:upper-oneway}) the cap $C_\varrho$ is the normalized
coefficient of its one-way repair \ref{it:record:upper-F7}
(\cref{lem:tr:repair}), with the joint-star estimate
\eqref{eq:frontier:all-prime-star-cap}, valid at every odd prime
(\cref{lem:upper-stars}(b), \ref{it:record:upper-F8}), on the ten $EL/LE$
records of degree two. The caps are evaluated at $2^{54}$
(\cref{cert:tail-caps}), where the spectral constants $g,B_{\rm top}$ of the
nonroot cuts that \ref{it:record:upper-F7} reads come from the tests of both
classes (\cref{cert:upper-band-inputs}(d)); by \cref{lem:record:domain}(iii)
the caps bound the repair coefficients at every prime $p\ge2^{54}$ of either
class. No class-three datum of the bands is used. With
$\Lambda^{\rm tail}_k(\mathfrak t)$ the sum of the caps of the repaired
records of state $\mathfrak t$ and degree $k$, put
\begin{equation}\label{eq:tail-polynomials}
 \Pi^{\rm rep}(r)=r^3\sum_{\mathfrak t}w_{\mathfrak t}\Lambda^{\rm tail}_2(\mathfrak t)
 +r^4\sum_{\mathfrak t}w_{\mathfrak t}\Lambda^{\rm tail}_3(\mathfrak t)
 =c^{\rm rep}_3r^3+c^{\rm rep}_4r^4,\qquad
 \Pi^{\rm del}(r)=\sum_{d=3}^7c^{\rm del}_dr^d,
\end{equation}
with $c^{\rm del}_d$ the total weight of the deleted records of overall
degree $d$, as in \cref{def:upper-comparison}. Every record of degree at
least four is deleted, so $c^{\rm del}_d=c_d$ for $4\le d\le7$, and
\[
 c^{\rm rep}_3=\tfrac{73056854152246473}{12500000000},\qquad
 c^{\rm rep}_4=\tfrac{149151332035400049}{68359375},\qquad
 c^{\rm del}_3=\tfrac{81606792}{25}.
\]

\begin{certificate}[The tail repair caps]\label{cert:tail-caps}
Each of the $94$ one-way repairs at a tail prime satisfies its
repair inequalities \ref{it:record:upper-G1},
\ref{it:record:upper-G2} and \ref{it:record:upper-G5} with the
coefficients evaluated at $2^{54}$ (hence at every larger prime,
by monotonicity), with the constants $g,B_{\rm top}$ of the
nonroot layer cuts at $2^{54}$
(\cref{cert:upper-band-inputs}(d)). The coefficients are stored in
\path{research/frontier_2026_10_17/sparse/high72_oneway_repairs.json}
(degree two) and \path{high72_oneway_repairs_degree3.json}
(degree three) in the same directory; the six $LR$ and four $RL$ records of degree
two are replaced by their joint-star rows \ref{it:record:upper-F8} at $2^{54}$
in \path{high72_oneway_jointstar.json}, whose hypotheses hold there. The
adapters \path{high72_oneway_adapter.py}
and \path{high72_oneway_jointstar_adapter.py} of that directory authenticate
the records against \path{high72_oneway_checks.json} and
\path{high72_oneway_star_checks.json}; no class-three datum of the bands is
used. The record of the tail datum,
\path{research/roadmap_2026_10_18/item12/upper_v2_local.json}, carries the
resulting ledgers $\Lambda^{\rm tail}_2,\Lambda^{\rm tail}_3$ and the
coefficients $c^{\rm rep}_3,c^{\rm rep}_4,c^{\rm del}_d$ above as its field
\texttt{tail\_comparison}; the independent review
\path{research/push_2026_10_24/tools/check_upper.py} compares them with the
tail data of the source it reviews and re-sums the tail prices.
\end{certificate}

\begin{lemma}[The tail comparison law and its prices]\label{lem:tail-prices}
Let $p\ge\Pfour$ be prime and $r=p^{-1/2}$. Let $\nu_p$ be obtained from
$\widetilde\mu_p$ by keeping every record that accepts the twelve queries,
replacing each of the $94$ nonaccepting two-layer records of degree two
and three by its one-way repair with the all-prime cap, deleting every other
nonaccepting record, and dividing by $P_\mu(r)$.
\begin{enumerate}[label=(\roman*)]
\item $\nu_p$ is a nonnegative square-affine invariant measure with the
 normalizer of $\mu_p$; every retained record accepts $\mathcal E_6$; the
 normalized accepted part has uniform one-slot marginals and query-pair
 marginals $\Unif(\Om)$; under every accepted part each query pair is a
 strict-square pair at $p$, which is \ref{it:record:R5} at the tail; and
 $\nu_p$ is admissible (\cref{def:admissible}).
\item $\nu_p-\mu_p$ satisfies the bracket of \cref{def:record:comparison-law}
 at $\iota$ with prices $s_p\le(\Pi^{\rm rep}+\Pi^{\rm del})(r)\,p^{2a}$ and
 $a_p\le\Pi^{\rm rep}(r)\,p^{2a}$.
\item $\eta_p:=s_p+a_p\le C_{\rm tail}\,p^{-3/2+2a}$, with
 $C_{\rm tail}=\sum_d(2c^{\rm rep}_d+c^{\rm del}_d)\Pfour^{-(d-3)/2}
 =1.4953\ldots\cdot10^7$ ($c^{\rm rep}_d=0$ for $d\ge5$).
\item With $w_d=d/2-2a$ and $\Phi_d(a)=\Pfour^{-w_d}+\Pfour^{1-w_d}/(w_d-1)$,
 for every $N$
 \[
  S_{\rm tail}:=\sum_{\Pfour\le p\le N^{1/20}}s_p\le
  \sum_{d=3}^7\bigl(c^{\rm rep}_d+c^{\rm del}_d\bigr)\Phi_d(a),\qquad
  A_{\rm tail}:=\sum_{\Pfour\le p\le N^{1/20}}a_p\le
  \sum_{d=3}^4c^{\rm rep}_d\Phi_d(a),
 \]
 and $\sum_{p\ge\Pfour}\eta_pp^{1/4}<\infty$. The tail atoms carry no
 first-stage and no bracket records, and their matching-envelope diagonals
 have price $T_{{\rm L},p}=T_{{\rm R},p}\le\eta_p$; so their part of the sum
 in (CP8)(c) is at most $3\sum_{p\ge\Pfour}\eta_pp^{1/4}<\infty$.
\end{enumerate}
\end{lemma}

\begin{proof}
(i), (ii). By \cref{lem:tr:repair}(a)--(c) a repaired record becomes a
probability law with exactly the actual half marginals at $\iota$ of the
record it replaces, supported in the support of the uniform law of the
repaired graph (half graphs kept, cross constraints removed, projected
queries imposed), so it accepts the twelve queries; its data are
square-affine invariant, so it is invariant. Its difference from the
replaced record has zero marginals and norm at most $C_\varrho r$ relative to
these marginals (\cref{lem:tr:repair}(f) with \ref{it:record:upper-F7}, the
cap bounding the coefficient at $p$ by \cref{lem:record:domain}(iii)). From
here the proof of \cref{prop:record:upper}(i), its admissibility paragraph
included, applies to the atom $\{p\}$ with $C^{\rm s}=C^{\rm a}=C_\varrho$.
Each repair bound is included in full in both the symmetric
total $s_p$ and the skew total $a_p$ (an overestimate in each),
which is allowed:
\cref{rem:cp:charging}(iii) would allow the skew part of the new law alone,
and the recorded tail polynomials do not use this. A repaired record of
state $\mathfrak t$ is charged at most $w_{\mathfrak t}C_\varrho
r^{\deg\mathfrak t+1}$ in each part, and summing gives $\Pi^{\rm rep}(r)$; a
deleted record is positive and self-adjoint at $\iota$, charged its mass in
$s_p$ only, and summing gives $\Pi^{\rm del}(r)$; the repaired records are
exactly those missing from $\Pi^{\rm del}$; repairs keep mass and deletions
decrease it, and dividing by $P_\mu(r)\ge1$ only decreases the prices, which
carry the factor $p^{2a}$ of an atom $\{p\}$ (\cref{def:price}). Every record
of $\nu_p$ is invariant, so the accepted part has uniform one-slot marginals
and invariant query-pair marginals supported on $\Om$, which are therefore
$\Unif(\Om)$. Every retained or repaired record imposes the strict-square
relation on each query pair.

(iii) Both polynomials start in degree three and $p^{-(d-3)/2}\le
\Pfour^{-(d-3)/2}$, so $\eta_p\le(2\Pi^{\rm rep}+\Pi^{\rm del})(r)p^{2a}\le
C_{\rm tail}p^{-3/2+2a}$.

(iv) For $w>1$ and an integer $P\ge2$, $\sum_{n\ge P}n^{-w}\le
P^{-w}+\int_P^\infty x^{-w}dx=P^{-w}+P^{1-w}/(w-1)$. Every $w_d\ge3/2-2a>1$,
the tail primes are integers $n\ge\Pfour$, and by (ii) the prices are sums
of $(c^{\rm rep}_d+c^{\rm del}_d)p^{-w_d}$, respectively $c^{\rm rep}_dp^{-w_d}$.
Next $\eta_pp^{1/4}\le C_{\rm tail}p^{-5/4+2a}$ and $-5/4+2a<-1$. Finally the
records of the first stage and of the bracket vanish outside the middle blocks
(\cref{def:record:frame}), so $p_p=v_p=d_p=u_p=0$ at the tail, and outside the
middle the
matching-envelope diagonals are the norm envelopes of the changes
(\cref{prop:record:matching}(i)), whose price is at most $s_p+a_p$.
\end{proof}

The upper prices $S_{\rm u}(a),A_{\rm u}(a)$ of the closing contain the two
sums of (iv) (\cref{prop:record:upper}(ii)). The final assembly adds them to
the termwise repriced band ledger in
\path{research/continuation_2026_10_24/global/closing.py} (\texttt{upper\_price},
lines 141--149). At the final exponent the tail contributes about
$2.97\cdot10^{-70}$ to $S$ and $1.90\cdot10^{-70}$ to $A$. The allowance
$10^{-12}$ of the closing covers the prefix exception and also a second,
repriced copy of the tail prices (\cref{cert:record:frame-data}(f)); the tail
is thus charged twice, which only increases $\Theta$.

\subsection{The external index by exact conductor}\label{subsec:index-law}

Each query has an interval coordinate modulo $K$ and independent external
prime-power coordinates. At $m=64$ the modulus used in the
square-class kernel (called \emph{physical} below, to
 distinguish it from the comparison-depth modulus) is
\[
 \Mzero=2^{129}\prod_{\substack{p<Q_3\\p\equiv3\ (4)}}p^{128}.
\]
The odd product includes $3$. At each prime choose, for each query, a uniform residue $t$ modulo the corresponding prime power, conditioned on $v_p(t)<m$, and use $t^2$ as the increment $r_{e,p}$ at $p$ (\cref{def:hyb:coupled}).
These increments are nonzero squares; their possible nonunit valuations
lie in the classes treated in \cref{lem:kernel-depth-classes}, applied in the proof of \cref{lem:imp:padic-envelope} below.
The random depth $J$ of \cref{lem:imp:padic-envelope} below
(the comparison depth) does not change this physical modulus. It
only determines the conditional expectation used in the positive
index envelope. Thus the cancellation in the actual kernel and
the comparison moment use different, explicitly specified moduli.

\begin{lemma}[Finite conductor-depth index envelope]
\label{lem:imp:padic-envelope}
Fix $m\ge3$ and $Q>2$. At every odd external prime
$p\equiv3\pmod4$ with $p<Q$, let $t$ be uniform modulo $p^{2m}$
conditioned on $p^m\nmid t$, and use $t^2$ as the increment (a
uniform conditioned root, squared).  At $2$ use the same
construction modulo $2^{2m+1}$.  These finite probability laws
have uniform endpoint marginals and preserve accepted-square cancellation, in that every increment lies in a class of \cref{lem:kernel-depth-classes}.  They admit the following decreasing conductor
tails, where $e_p=p^{-m}$:
\[
 \begin{array}{ll}
 p\text{ odd}:&\delta_0=1,\quad
 \delta_k=\dfrac{p^{-k/2}+e_p}{1-e_p}\ (1\le k<2m),
 \quad\delta_{2m}=0;\\[4pt]
 p=2:&\delta_0=1,\quad
 \delta_1=\delta_2=\dfrac{2^{-1/2}+e_2}{1-e_2},\\[4pt]
 &\delta_k=\dfrac{2^{-(k-1)/2}+e_2}{1-e_2}\ (3\le k<2m+1),
 \quad\delta_{2m+1}=\dfrac{e_2}{1-e_2},\quad
 \delta_{2m+2}=0.
 \end{array}
\]
Independent depths $J$ with $\Pr(J\ge k)=\delta_k$ therefore supply
the finite-filtration index envelope of \cref{lem:record:common-mean},
and its conclusion (a) retains the common mean.
\end{lemma}
\begin{proof}
For odd $p$, the primitive quadratic Gauss multiplier at conductor
$p^k$ satisfies $G_k=p^{-1}G_{k-2}$, $G_0=1$ and
$|G_1|=p^{-1/2}$.  Summing first over the last base-$p$ digit
proves the recurrence.  Excluded roots have square $0$ modulo
$p^{2m}$, so the retained multiplier is
$(G_k-e_p)/(1-e_p)$.  At the even top conductor $G_{2m}=e_p$
and the multiplier is exactly zero.  At $2$, $G_1=0$ and
$|G_k|=2^{-(k-1)/2}$ for $k\ge2$, by \cref{thm:gauss} applied to the quadratic Gauss sum
modulo $2^k$.  The excluded roots
have square $0$ below the top conductor, while their top phases
cancel in equal numbers.  This proves all the displayed bounds.
For $m\ge3$ the $\delta_k$ are nonincreasing in $k$ and $\delta_1<1$, so they define probability depth laws.

\emph{The comparison envelope.} Take $E_j$ to be conditional expectation
on the residue modulo $p^j$. Its exact-depth space $E_j-E_{j-1}$
consists of the characters of conductor $p^j$. The Gauss-sum bounds
just proved therefore give the operator-norm bounds on these
spaces required by \cref{lem:record:common-mean}. Independence of the
prime coordinates gives the product envelope. The field Gram commutes
with the residue projections, and the constant multiplier is one, so
that lemma gives both signed-index bounds and retains the common mean;
self-adjointness of the transition is not required.

\emph{Physical cancellation.} Every retained root has valuation $j<m$.
Its square modulo $p^{2m}$ has valuation $2j<2m$ for odd $p$; at $2$
its square modulo $2^{2m+1}$ has $2j+3\le2m+1$. These are exactly the
nonvanishing square classes of \cref{lem:kernel-depth-classes}.
That lemma supplies the normalized periodic root weights, their uniform
higher-digit distribution and their Fourier bounds. Thus the
actual-square kernel applies to every increment, with the full physical
modulus fixed; no increment is $0$ modulo $p^{2m}$ (resp.\ $2^{2m+1}$ at $2$).

The full physical modulus remains pinned in this cancellation,
independently of the random comparison depth.  In the comparison,
conditioning the normalized input gives exactly
\[
 \mathbb E[p^H\mathbf1_{n=s\ (p^H)}\mid s\ (p^j)]
        =p^j\mathbf1_{n=s\ (p^j)}.
\]
The exact-fibre and square-step arguments of \cref{lem:record:fibre-scale}(a), (b) hence apply without change.  All moduli are fixed before $N$ varies, so
only constants independent of $N$ and the eventual starting scale increase.
Uniform averaging of the input over residues and intervals is
still the original lift.  This retains the density-linear mass
and all common means.
\end{proof}

\begin{proposition}[Finite p-adic matching and wrap prices]
\label{prop:imp:padic-price}\label{prop:oct24-index-cost}
With $a$ the price exponent and $m,Q$ as in
\cref{lem:imp:padic-envelope}, and $B,S,A,T$ as in
\cref{thm:comparison}, put $z_p=p^{a/18}$. The exact one-query
conductor moment is
\[
 M_{\rm pad}(K)=K^{a/36}L_{2,m}
       \prod_{\substack{p<Q\\p\equiv3\ (4)}}L_{p,m},
\]
where, with $e=p^{-m}$,
\begin{align*}
 L_{p,m}&=1+\frac{(z_p-1)p^{-1/2}
                    [1-(z_p/p)^m]/(1-z_p/p)
                       +e(z_p^m-1)}{1-e}\quad(p\text{ odd}),\\
 L_{2,m}&=1+\frac{(z_2-1)
           [2^{-1/2}+z_2(1-(z_2/2)^{m-1})/(2-z_2)]
                  +2^{-m}(z_2^{m+1}-1)}{1-2^{-m}}.
\end{align*}
For every fixed integer $K\ge2$, the matching and wrap prices are
\begin{equation}\label{eq:oct24-index-cost}
 J=Be^T\bigl(M_{\rm pad}(K)/2\bigr)^6,\qquad
 W=Be^{S+A}\left[\left(1+\frac{M_{\rm pad}(K)}
                                  {2K^{35/36}}\right)^6-1\right].
\end{equation}
The coefficient $63/64=1-2^{-6}$ in \cref{thm:comparison} and all field restoration terms in the closing criterion there remain unchanged.
\end{proposition}
\begin{proof}
For each depth $J$, the least square multiple of $p^J$ (the
smallest square divisible by $p^J$) is $p^{2\lceil J/2\rceil}$.  Its two occurrences among the
$72$ slots give exponent $a/36$.  Summation by parts therefore
gives
\[
 L_{p,m}=1+(z_p-1)\sum_{l\ge1}z_p^{l-1}\delta_{2l-1}.
\]
Substituting the tails in \cref{lem:imp:padic-envelope} gives
the displayed formulas.  In particular the top 2-adic depth
has $\lceil(2m+1)/2\rceil=m+1$ and is included.  The entire $2$-part and prime $3$ are represented by their depth
factors; no additional fixed conductor factor is inserted.

Fix the depths, apply the folds $\phi_\pm$ of \cref{def:sorting}, then average over the depths.
There are exactly two occurrences per query, and six independent
queries.  \Cref{lem:record:direct-gram} gives $J$.  For a nonempty
subset of $j$ active queries in the signed expansion, any
nonwrapping interval permits accepted-square cancellation.  All
$j$ wrap with probability $K^{-j}$.  The density normalization multiplies their bound by $K^{j/36}$ and their depth average is $M_{\rm pad}(K)^j$.
The unchanged positive field envelope is $Be^{S+A}$.  Summing
the coefficients $\binom6j2^{6-j}/64$ gives $W$.  The common
mean is unchanged by uniform index averaging, as proved above.
\end{proof}

\begin{corollary}[A compact upper majorant]
\label{cor:imp:padic-majorant}
Take $m=64$, $0<a\le3/2500$, $2^{16}\le Q\le2^{40}$, and
put $Z=Q^{a/18}$.  Define
\[
 L_{p,\infty}=1+\frac{\sqrt p(z_p-1)}{p-z_p},\qquad
 L_{2,\infty}=1+(z_2-1)\left(2^{-1/2}+\frac{z_2}{2-z_2}\right),
\]
\[
 E_o=\frac{Q3^{-m}Z^m}{1-3^{-m}},\qquad
 E_2=\frac{2^{-m}z_2^{m+1}}{1-2^{-m}}.
\]
Let $M_{\rm old}^{+}$ be any proved upper bound for $M_a(K)$ of \eqref{eq:oct24-moment} below, and set
$F_p=1+\sqrt{p+1}(z_p-1)/(p-1)$.  For any exact prime cutoff
$7\le U\le Q$, a valid upper bound for $\log M_{\rm pad}(K)$ is
\begin{align*}
 \log M_{\rm old}^{+}-\frac a{36}\log144
  +\log L_{2,\infty}+\log L_{3,\infty}
  -\sum_{\substack{7\le p<U\\p\equiv3\ (4)}}
       \log\frac{F_p}{L_{p,\infty}}+E_o+E_2.
\end{align*}
The error term satisfies $E_o+E_2<2^{-60}$ on this domain.
\end{corollary}
\begin{proof}
Discarding the negative geometric tails in the exact finite
formulas gives
\[
 L_{p,m}\le L_{p,\infty}
       \left(1+\frac{p^{-m}z_p^m}{1-p^{-m}}\right),\qquad
 L_{2,m}\le L_{2,\infty}
       \left(1+\frac{2^{-m}z_2^{m+1}}{1-2^{-m}}\right).
\]
There are fewer than $Q$ odd primes, so $\log(1+x)\le x$
gives the two surcharges.  Also
$Z\le2^{1/375}<501/500<5/4$.  For $p\ge3$, $1\le z\le5/4$,
\[
 (p+1)(p-z)^2-p(p-1)^2
 \ge\frac{8p^2-31p+25}{16}\ge\frac14.
\]
Thus $L_{p,\infty}\le F_p$.  Replace only the explicitly
enumerated factors by their smaller values and keep the remaining
old upper bound.  Replacing its factor $144$ by the two new factors
proves the claimed logarithmic majorant.  Finally the rational
inequality $(501/500)^{375}>2$ and direct substitution at
$Q=2^{40}$ give
\[
 \frac{2^{40}3^{-64}(501/500)^{64}}{1-3^{-64}}
 +\frac{2^{-64}(501/500)^{65}}{1-2^{-64}}<2^{-60}.
\]
$L_{p,\infty}$ and $L_{2,\infty}$ are used here only as majorants, with the displayed surcharges, for the finite-depth factors $L_{p,m}$ and $L_{2,m}$.  An interval
enclosing this upper majorant supplies no lower bound for the
physical moment.
\end{proof}


\subsection{Evaluating \texorpdfstring{$R_I$}{R\_I}, \texorpdfstring{$J$}{J} and \texorpdfstring{$W$}{W}}

The remaining quantities in $\Theta$ coming from the external
index are its conductor moment $R_I$, and through it $J$ and $W$.
The moment of the index law used in the construction is
$M_{\rm pad}(K)$ of \cref{prop:imp:padic-price}; its evaluation
uses a simpler auxiliary moment, a product of elementary factors
over the $634017$ primes $p\equiv3\pmod4$ below $Q_3$:

\begin{equation}\label{eq:oct24-moment}
 M_a(K)=(144K)^{a/36}
 \prod_{\substack{7\le p<Q_3\\p\equiv3\ (4)}}
 \left[1+\frac{\sqrt{p+1}}{p-1}(p^{a/18}-1)\right].
\end{equation}
The following lemma encloses this product using only $22$
explicit prime sums, so that it can be certified without
multiplying out $634017$ factors in interval arithmetic; this
supplies the upper bound used in \cref{cor:imp:padic-majorant}.

\begin{lemma}[Finite prime moments]\label{lem:index-moments}
Let $Q>7$ be an integer and $a\ge0$; sums and products over $p$ run over the
primes $p\equiv3\pmod4$ with $7\le p<Q$, and $\delta_p=\sqrt{p+1}/(p-1)$.
For $1\le j\le4$ and $j\le k\le7$ put
\[
 M_{jk}(Q)=\sum_p\delta_p^j\Bigl(\frac{\log p}{18}\Bigr)^k,\qquad
 P_j(a)=\sum_p\bigl[\delta_p\bigl(p^{a/18}-1\bigr)\bigr]^j,
\]
and
\[
 T_j(a)=\sum_{k=j}^6\frac{j!}{k!}\genfrac{\{}{\}}{0pt}{}{k}{j}a^kM_{jk}(Q),
 \qquad
 \varrho_j(a)=e^{ja\log Q/18}\,\frac{j^7a^7}{7!}\,M_{j7}(Q),
\]
with the Stirling numbers $\genfrac{\{}{\}}{0pt}{1}{k}{j}$ of the second kind.
Then $T_j(a)\le P_j(a)\le T_j(a)+\varrho_j(a)$, and
$G_Q(a)=\prod_p\bigl(1+\delta_p(p^{a/18}-1)\bigr)$ satisfies, with every term
evaluated at $a$,
\[
 T_1-\tfrac12(T_2+\varrho_2)+\tfrac13T_3-\tfrac14(T_4+\varrho_4)\;\le\;\log G_Q(a)
 \;\le\;T_1+\varrho_1-\tfrac12T_2+\tfrac13(T_3+\varrho_3).
\]
In particular, for every integer $\Kint\ge2$, the moment
$M_a(\Kint)=(144\Kint)^{a/36}G_{Q_3}(a)$ of \eqref{eq:oct24-moment} is enclosed by the explicit bounds above, as a function of $a$, $\Kint$, $Q_3$ and the $22$ positive sums
$M_{jk}(Q_3)$.
\end{lemma}

\begin{proof}
\emph{Taylor bounds.} For $y\ge0$ the exponential generating function of
surjections gives
$(e^y-1)^j=j!\sum_{k\ge j}\genfrac{\{}{\}}{0pt}{1}{k}{j}y^k/k!$, a series with
nonnegative coefficients; its terms of degree at most six are a lower bound.
Since $j!\genfrac{\{}{\}}{0pt}{1}{k}{j}$ counts the surjections of a $k$-set
onto a $j$-set, it is at most $j^k$, the number of all maps; so the terms of
degree at least seven are at most
$\sum_{k\ge7}(jy)^k/k!\le(jy)^7e^{jy}/7!$. Take $y=a(\log p)/18$, which is at
most $a\log Q/18$ for $p<Q$, multiply by $\delta_p^j$ and sum over $p$: as
$p^{a/18}-1=e^y-1$, this gives $T_j(a)\le P_j(a)\le T_j(a)+\varrho_j(a)$.

\emph{Logarithm.} For $x\ge0$ and $0\le s\le x$,
$1-s+s^2-s^3\le1/(1+s)\le1-s+s^2$, since $1/(1+s)=1-s+s^2-s^3/(1+s)$;
integrating over $[0,x]$ gives
$x-\frac{x^2}2+\frac{x^3}3-\frac{x^4}4\le\log(1+x)\le x-\frac{x^2}2+\frac{x^3}3$.
Apply this at $x=\delta_p(p^{a/18}-1)\ge0$ and sum over $p$:
$P_1-\frac12P_2+\frac13P_3-\frac14P_4\le\log G_Q(a)\le P_1-\frac12P_2+\frac13P_3$.
Inserting the bounds of the first step, each with the sign of its
coefficient, gives the two displayed bounds.

\emph{The moment.} By \eqref{eq:oct24-moment}, $M_a(\Kint)=(144\Kint)^{a/36}G_{Q_3}(a)$,
and the bounds above involve, besides $a$, $\Kint$ and $Q_3$, only the sums
$M_{jk}(Q_3)$, $1\le j\le4$, $j\le k\le7$. These are finite sums of positive
terms, so their enclosures need only outward square roots and logarithms of
the primes. This auxiliary product is the Bernoulli moment from
\cref{lem:hyb:moment,lem:square-shift}; it prices the remaining large-prime
factors in \cref{cor:imp:padic-majorant}, rather than defining the
selected physical index law.
\end{proof}

\begin{remark}[The choice of $K$]\label{rem:index-convexity}
For a product majorant $M_{\rm pad}(K)\le\bar c K^{a/36}$ with $M_{\rm pad}$ as in \cref{prop:imp:padic-price}, and $0<a<35$, the bound there for
$J+W$ is a positive sum of $e^{ax/6}$ and
$e^{-j(35-a)x/36}$, $1\le j\le6$, where $x=\log K$.
It is convex. The baseline integer is $K=\PushK$; its validity
requires only evaluation of the two prices, not a proof of optimality.
\end{remark}

\begin{certificate}[The index moment]\label{cert:index-moment}
The auxiliary sums $M_{jk}(Q_3)$ run over the $634017$ primes
$p\equiv3\pmod4$ with $7\le p<19951586$. The largest is $19951579$
and their sum is $6108518317099$. Their outward 48-digit enclosures
and the prime-list digest are retained in
\path{research/continuation_2026_10_24/global/prime_moments.json}.
The checker of \cref{app:certificates} regenerates these finite prime data.

For the baseline finite law, use $m=64$, $Q_3=19951586$,
$a=\PushA$, $K=\PushK$ and small-prime replacement cutoff
$U=65536$ in \cref{cor:imp:padic-majorant}. The outward evaluation
and its independent reconstruction are
\path{research/improvements_2026_10_25/global/padic_moment.py} and
\path{research/improvements_2026_10_25/global/check_padic_moment.py}.
The selected record is
\path{research/push_2026_10_25/artifacts/padic_moment.json}.
They certify
\[
 R_I=M_{\rm pad}(K)\le
 1.247139907406884705556794769493.
\]
This rational is the upper bound used for $R_I$ in \cref{def:closing-quantities}(12). The narrow recorded interval encloses an upper majorant;
its lower endpoint is not a lower bound for $M_{\rm pad}(K)$.
Both scalar assemblies, author and independent, use the same upper rational,
and \path{research/push_2026_10_25/verify.py} reruns its independent
check (\cref{cert:closing}).
\end{certificate}

\paragraph{Consequences and what to read next.}
The tail datum (\cref{def:record:upper-bands}(e)),
\cref{def:upper-reference,def:upper-comparison},
\cref{subsec:layer72:upper,subsec:upper-reference-rp}, \cref{lem:upper-geometry},
\cref{lem:upper-reference-rp,cert:upper-completion},
\ref{it:record:upper-F1}--\ref{it:record:upper-F8}, with
\ref{it:record:upper-F5}(b), and \cref{lem:record:domain}, the class-one
inputs (\cref{subsec:upper-nonroot,subsec:advance:product-rows},
\cref{cert:upper-band-inputs}(d),
\cref{lem:upper-completed-root,lem:spec2:paley-one,lem:spec2:parity-transfer}),
\cref{sec:frontier:upper-oneway}, \cref{lem:tr:repair,lem:upper-stars} with
the caps of \cref{cert:tail-caps}, \cref{rem:cp:charging},
\cref{prop:record:upper,prop:record:matching},
\cref{lem:trace-mixing,lem:mixture}, (S7) of \cref{cert:record:analytic-stage},
\cref{lem:square-shift}, \cref{def:hyb:coupled,lem:hyb:moment},
\cref{lem:imp:padic-envelope,prop:imp:padic-price,cor:imp:padic-majorant,lem:index-moments}
and \cref{thm:comparison}.

\paragraph{Hypotheses verified in this section.}
(CP1) for the tail atoms of both classes (\cref{lem:tail-law});
\ref{it:record:R4} (\cref{lem:record:tail-mixing}); \ref{it:record:R5}, the
tail part of (CP4) and the admissibility (CP8)(a) of the tail laws
(\cref{lem:tail-law}(ii), \cref{lem:tail-prices}(i),(ii)); the summability (CP8)(c)
at the tail (\cref{lem:tail-prices}(iii),(iv)); and (CP0)(d) at
$\Kint=\PushK$ (\cref{prop:oct24-index-cost,lem:index-moments,cert:index-moment}).
These are the rows ``(CP1) tail'', ``(R4)'', ``(R5), (CP4)(a)'' (at the tail),
``(CP4) tail'', ``(CP8)(c)'' and ``(CP0)(d) index law'' of
\cref{tab:hypotheses}.
